% /modeling/sparse.tex

\chapter{Sparse Alignment Model}

\section*{Motivation}

% modeling/sparse.tex
\begin{figure}[htbp]
\centering
% figures/modeling/sparse_model_structure.tex
\begin{tikzpicture}[
  node distance=0.65cm,
  elem/.style={draw, rounded corners, align=center, minimum width=2.25cm, minimum height=0.9cm},
  arrow/.style={->, thick}
]
\node[elem, fill=blue!8] (f1) {Fixed\\\small $G/R$};
\node[elem, fill=orange!12, right=of f1] (t1) {Transition\\\small $\kappa(s)$};
\node[elem, fill=blue!8, right=of t1] (f2) {Fixed\\\small $G/R$};
\node[elem, fill=orange!12, right=of f2] (t2) {Transition\\\small $\kappa(s)$};
\node[elem, fill=blue!8, right=of t2] (f3) {Fixed\\\small $G/R$};

\draw[arrow] (f1) -- (t1);
\draw[arrow] (t1) -- (f2);
\draw[arrow] (f2) -- (t2);
\draw[arrow] (t2) -- (f3);

\node[draw, rounded corners, fill=green!10, align=center, below=1.2cm of f2, minimum width=6.2cm]
{sparse alignment = compact curvature program};
\end{tikzpicture}

\caption{Sparse alignment as an alternating sequence of fixed and transition elements.}
\label{fig:sparse-model-structure}
\end{figure}

\begin{remark}
A continuous curvature law is mathematically natural, but it is not by
itself an operational representation for editing, exchange, validation,
or reconstruction from heterogeneous engineering data.
For this reason, a sparse model is introduced as a structured,
piecewise-defined representation of an alignment.
\end{remark}

\begin{remark}
The sparse model is not intended as a denial of continuous geometry.
On the contrary, it serves as an intermediate layer between continuous
curvature theory and practical modelling tasks such as import,
computation, optimization, and interoperability.
\end{remark}

\begin{remark}
Dense geometric representations such as point clouds or sampled
polylines preserve positional information, but largely destroy the
structural semantics of the alignment itself.
\end{remark}

% ============================================================
\section{General Idea}
% ============================================================

\begin{definition}[Sparse alignment model]
A sparse alignment model is a structured longitudinal representation
consisting of synchronized alignment bands together with sufficient
reconstruction information.
\end{definition}

\begin{remark}
The adjective \emph{sparse} refers here to the fact that the alignment
is not stored as a dense point cloud, but as a compact sequence of
semantically meaningful elements.
\end{remark}

\begin{remark}
Within AIM, the sparse model is not interpreted as a static 3D geometry,
but as a persistent structural kernel from which geometry and dynamics
are derived.
\end{remark}

% ============================================================
\section{Element Types}
% ============================================================

\begin{definition}[Alignment element]
An alignment element is a segment of an alignment with a prescribed
local curvature behaviour on a parameter interval.
\end{definition}

\begin{definition}[Fixed element]
A fixed element is an alignment element with constant curvature
\[
\kappaf(s)=\kappaf_0.
\]
Typical special cases are straight segments with $\kappaf_0=0$ and
circular arcs with $\kappaf_0\neq 0$.
\end{definition}

\begin{definition}[Transition element]
A transition element is an alignment element with non-constant
curvature, i.e.
\[
\kappaf(s)=f(s)
\]
for a suitable curvature law $f$.
\end{definition}

\begin{remark}
Within AIM, transition elements are primarily interpreted through their
curvature evolution rather than through explicit geometric shape.
\end{remark}

\begin{remark}
At this level of abstraction, no commitment is yet made to a specific
family of transition laws.
This includes classical clothoids as well as polynomial, lookup-based,
or hybrid constructions.
\end{remark}

\TODO{Formal classification of transition families.}

% ============================================================
\section{Minimal Parameterization}
% ============================================================

\begin{definition}[Arc length of an element]
Each element $e_i$ is associated with an arc length
\[
L_i > 0.
\]
\end{definition}

\begin{definition}[Local curvature law]
Each element $e_i$ is associated with a curvature law
\[
\kappaf_i : [0,L_i] \to \R.
\]
\end{definition}

\begin{remark}
For fixed elements, $\kappaf_i$ is constant.
For transition elements, $\kappaf_i$ varies along the local parameter.
\end{remark}

\begin{definition}[Sparse alignment]
A sparse alignment is an ordered tuple
\[
A=(e_1,e_2,\dots,e_n)
\]
together with sufficient initial data for geometric reconstruction.
\end{definition}

% ============================================================
\section{Initial Data and Pose}
% ============================================================

\begin{definition}[Planar start pose]
A planar start pose is a tuple
\[
P_0=(x_0,y_0,\theta_0)
\]
with position $(x_0,y_0)\in\R^2$ and initial direction
$\theta_0\in\R$.
\end{definition}

\begin{proposition}
The planar alignment is reconstructed by successive integration of the
curvature law:
\[
\kappa(s)
\rightarrow
\theta(s)
\rightarrow
\gamma(s).
\]
\end{proposition}

\begin{remark}
Together with the ordered element sequence and the local curvature laws,
the start pose determines the alignment geometry by successive
integration.
\end{remark}

\begin{definition}[Reconstructable sparse alignment]
A sparse alignment is called reconstructable if its element sequence and
its initial data suffice to determine a unique planar alignment up to
the intended modelling conventions.
\end{definition}

% ============================================================
\section{Sparse Longitudinal Band Structure}
% ============================================================

\begin{principle}[Sparse longitudinal band principle]
Railway alignments are represented as synchronized longitudinal bands
rather than as monolithic geometric objects.
\end{principle}

\begin{remark}
Within the AIM framework, the persistent alignment model is intentionally
kept sparse and decomposed into synchronized longitudinal bands.
\end{remark}

% ------------------------------------------------------------
\subsection{cGeom}
% ------------------------------------------------------------

\begin{definition}[Core geometric band]
The central geometric band is denoted by
\[
\mathrm{cGeom}.
\]
It represents the curvature-driven planar alignment core and consists of:
\begin{itemize}
\item fixed elements,
\item transition elements,
\item pose2-based reconstruction anchors.
\end{itemize}
\end{definition}

\begin{remark}
The cGeom band forms the primary integration core of the alignment.
\end{remark}

% ------------------------------------------------------------
\subsection{Cant Band}
% ------------------------------------------------------------

\begin{definition}[Cant band]
Cant is represented as an independent longitudinal band correlated with
fixed alignment elements.
\end{definition}

\begin{remark}
The cant band:
\begin{itemize}
\item does not define an independent stationing system,
\item stores rail height differences rather than roll angles,
\item follows the same normalized curvature family as the associated
cGeom element,
\item may contain local station offsets at element boundaries.
\end{itemize}
\end{remark}

\begin{remark}
Since cant evolution follows the same normalized curvature family as
the associated cGeom element, no independent geometric evolution model
is required for the cant band itself.
\end{remark}

\begin{remark}
The corresponding roll angle is derived via:
\[
\phi(s)
=
\arctan\!\left(
\frac{u(s)}{g}
\right),
\]
where:
\begin{itemize}
\item \(u(s)\): rail height difference,
\item \(g\): gauge.
\end{itemize}
\end{remark}

% ------------------------------------------------------------
\subsection{Profile Band}
% ------------------------------------------------------------

\begin{definition}[Profile band]
Vertical geometry is represented independently through a profile band.
\end{definition}

\begin{remark}
Thus, cGeom, cant, and profile are treated as synchronized but distinct
bands sharing a common arc-length domain.
\end{remark}

% ------------------------------------------------------------
\subsection{Runtime State}
% ------------------------------------------------------------

\begin{proposition}
Pose3 is not treated as persistent truth within the sparse model.
\end{proposition}

\begin{remark}
Persisting pose3 directly would introduce redundancy and synchronization
instabilities between geometry, cant, and profile representations.
\end{remark}

\begin{remark}
Instead, pose3 is derived dynamically from:
\begin{itemize}
\item cGeom,
\item cant,
\item profile,
\item and runtime evaluation services.
\end{itemize}
\end{remark}

\begin{remark}
Consequently, the primary intelligence of the AIM framework resides not
in the storage format itself, but in the mathematical operators and
evaluation services acting on the sparse model.
\end{remark}

% ============================================================
\section{Concatenation}
% ============================================================

\begin{definition}[Element concatenation]
Two consecutive elements are concatenated if the end state of the first
element provides the start state of the second element.
\end{definition}

\begin{definition}[Positional continuity]
A sparse alignment is positionally continuous if consecutive elements
meet in a common endpoint.
\end{definition}

\begin{definition}[Directional continuity]
A sparse alignment is directionally continuous if consecutive elements
share the same tangent direction at their common boundary.
\end{definition}

\begin{remark}
In most engineering contexts, positional and directional continuity are
minimal requirements. Higher continuity requirements may depend on the
chosen transition families and on the intended application.
\end{remark}

\TODO{Formalize continuity classes such as $C^0$, $C^1$, and curvature-related continuity conditions.}

% ============================================================
\section{Normal Form and Alternation}
% ============================================================

\begin{remark}
In practical sparse representations, a frequent pattern is the
alternation between fixed elements and transition elements. This
corresponds closely to engineering intuition and to many traditional
alignment descriptions.
\end{remark}

\OPEN{Should strict alternation between fixed and transition elements be formulated as a mathematical property, as a preferred normal form, or merely as an implementation convention?}

\begin{remark}
It is plausible to regard alternation as a normal form rather than a
fundamental axiom. Multiple neighbouring fixed elements may be merged,
and certain degenerate transition cases may collapse into fixed
elements.
\end{remark}

% ============================================================
\section{Normalized Representation of Elements}
% ============================================================

\begin{definition}[Normalized local parameter]
A normalized local parameter is a variable
\[
u\in[0,1]
\]
used to represent an element independently of its physical length.
\end{definition}

\begin{definition}[Normalized curvature law]
A normalized curvature law is a law defined on the unit interval,
typically of the form
\[
\widehat{\kappaf} : [0,1]\to\R.
\]
\end{definition}

\begin{remark}
Normalization separates shape from scale.
It allows the same transition family to be reused for different lengths
and different curvature magnitudes.
\end{remark}

\begin{remark}
Normalized representations form the basis for reusable transition
registries, family classification, and implementation-independent
comparison of alignment elements.
\end{remark}

% ============================================================
\section{Structural Role of the Sparse Model}
% ============================================================

\begin{remark}
The sparse model occupies an intermediate conceptual level:
\begin{itemize}
  \item below the purely abstract curvature-theoretic foundation,
  \item but above dense numerical sampling, imported point lists, or
        format-specific container structures.
\end{itemize}
\end{remark}

\begin{remark}
The sparse model is intentionally not a dense geometric storage format,
nor a direct copy of external exchange containers.
\end{remark}

\begin{remark}
This intermediate status makes the sparse model suitable as
\begin{itemize}
  \item an internal modelling layer,
  \item a reconstruction target from imported data,
  \item a source for numerical evaluation,
  \item and a bridge toward standards and BIM-related representations.
\end{itemize}
\end{remark}

% ============================================================
\section{Relation to Imported and Standardized Data}
% ============================================================

\begin{remark}
Existing engineering data rarely arrives directly in sparse form.
Instead, it is typically given through container formats, survey data,
legacy descriptions, or derived coordinate sequences.
\end{remark}

\begin{remark}
Such data may contain:
\begin{itemize}
\item inconsistent stationing,
\item incomplete cant information,
\item broken coordinate reference systems,
\item discontinuities,
\item partial geometry fragments,
\item or non-synchronized alignment components.
\end{itemize}
\end{remark}

\begin{remark}
The sparse model therefore should not be identified with any one
exchange standard. It is more appropriate to view it as a compact,
internal representation that can be derived from, and mapped back to,
multiple external formats.
\end{remark}

\RESEARCH{Dock sparse model to \railML\ alignment structures.}
\RESEARCH{Dock sparse model to \LandXML\ alignment representations.}
\OPEN{Define a minimal common abstraction layer for standards mapping.}
\OPEN{Clarify the relation between topology-oriented standards and purely geometric container models.}

% ============================================================
\section{Relation to BIM and BIMalignment}
% ============================================================

\begin{remark}
A sparse alignment model may support BIM-oriented workflows, but it
should not merely imitate BIM container logic.
Its primary responsibility is to represent alignment structure in a
mathematically and operationally coherent way.
\end{remark}

\RESEARCH{Formulate a BIM-compatible alignment core without losing geometric precision.}
\OPEN{Position the sparse model between generic BIM, infrastructure BIM, and explicit BIMalignment concepts.}
\TODO{State explicitly where the manuscript aims to push BIMalignment rather than only satisfy BIM conformity.}

% ============================================================
\section{Summary}
% ============================================================

\begin{summary}
The sparse alignment model is a compact, structured representation of an
alignment by means of synchronized longitudinal bands together with
reconstruction data and concatenation rules.

Within AIM, the sparse alignment model is not interpreted as a static
3D geometry, but as a persistent structural kernel from which geometry,
runtime state, and dynamic interpretation are derived.

It serves as a bridge between continuous curvature theory and practical
engineering operations such as import, validation, editing, numerical
evaluation, optimization, and interoperability.
\end{summary}
